% =====================================================================
% Chapter 7 --- Synthetic Dataset and 3D Detection Pipeline (AirSim)
%                                                       [detection -- OTHER PC]
% Self-contained project: Overview -> Dataset -> Detection methodology ->
% Results -> Discussion. This is the chapter completed on the other computer.
% Target length: ~9 pages (methodology + results).
% =====================================================================

\section{Overview and Objectives}
% TODO: motivation -- scarcity of annotated aerial 3D datasets; use AirSim to generate
% data and train the detection front-end of the pipeline (Chapter~\ref{cha:pipeline}).

\section{The AirSim Synthetic Dataset}

\subsection{Simulation Environment and Scenarios}
% TODO: AirSim/Unreal; urban scenarios; rationale for synthetic data and generalization.

\subsection{Sensor Channels and Annotations}
% TODO: RGB + depth + LiDAR + ego-pose; 2D/3D boxes, class, tracking IDs.

\subsection{Dataset Splits and Format}
% TODO: train/val/test; KITTI/nuScenes-compatible format.

\section{Detection Methodology}

\subsection{2D Detection}
% TODO: YOLO detector; inputs/outputs; training.

\subsection{Frustum Proposal}
% TODO: lifting 2D boxes into 3D frustums with camera intrinsics + point cloud.

\subsection{3D Detection with PointNet}
% TODO: PointNet on the frustum points; amodal 3D bounding box.

\section{Results}

\subsection{2D Detection Performance}
% TODO

\subsection{3D Detection Performance}
% TODO

\subsection{Generalization across Scenarios}
% TODO

\section{Discussion}
% TODO: how the trained detector feeds the tracker; sim-to-real outlook
% (link to the real-data tracking of Chapter~\ref{cha:monocular}).
